\documentclass[11pt]{article}

\usepackage[utf8]{inputenc}
\usepackage[T1]{fontenc}
\usepackage{lmodern}
\usepackage{geometry}
\usepackage{amsmath,amssymb,amsfonts,amsthm,mathtools}
\usepackage{bm}
\usepackage{enumitem}
\usepackage{microtype}
\usepackage{hyperref}
\usepackage{titlesec}
\usepackage{booktabs}
\usepackage{array}
\usepackage{setspace}
\usepackage{mathrsfs}
\usepackage{physics}

\geometry{margin=1in}
\hypersetup{colorlinks=true, linkcolor=black, urlcolor=blue, citecolor=black}
\setlist[enumerate]{topsep=0.4em,itemsep=0.35em}
\setlist[itemize]{topsep=0.3em,itemsep=0.25em}
\titleformat{\section}{\large\bfseries}{\thesection.}{0.5em}{}
\titleformat{\subsection}{\normalsize\bfseries}{\thesubsection.}{0.5em}{}
\setstretch{1.06}

\newcommand{\Aether}{\AE{}ther}
\newcommand{\Aflow}{\AE{}ther-flow}
\newcommand{\TheoryName}{\Aflow{} Interpretation of Relativity}
\newcommand{\TheoryProgram}{\Aether{} / \Aflow{} framework}
\newcommand{\ReBridge}{g^{\mathrm{br}}}
\newcommand{\OpMetricIER}{g^{\mathrm{IER}}}
\newcommand{\SigmaIER}{\Sigma^{\mathrm{IER}}}
\newcommand{\SigmaOff}{\Sigma^{\mathrm{IER,off}}}
\newcommand{\SigmaLongThree}{\Sigma^{\mathrm{IER},(3)}}
\newcommand{\SigmaC}{\Sigma^{\mathrm{C}}}
\newcommand{\SigmaReg}{\Sigma^{\mathrm{reg}}}
\newcommand{\ReOff}{\widetilde{\mathcal R}_{\mu\nu}^{\mathrm{off}}}
\newcommand{\ReOffTwo}{\widetilde{\mathcal R}_{\mu\nu}^{\mathrm{off},(2)}}
\newcommand{\ReLongThree}{\widetilde{\mathcal R}_{\mu\nu}^{(\chi,3)}}
\newcommand{\ReLongFour}{\widetilde{\mathcal R}_{\mu\nu}^{(\chi,\ge 4)}}
\newcommand{\DeltaTIER}{\Delta T_{\mu\nu}^{\mathrm{IER}}}
\newcommand{\EinLin}{\mathcal E_{\ReBridge}}
\newcommand{\EinQuad}{\mathcal Q_{\ge 2}^{\mathrm{Ein}}}
\newcommand{\EinInv}{\mathfrak G_{\mathrm{Ein}}}
\newcommand{\EinBil}{\mathcal B_{\ReBridge}^{\mathrm{Ein}}}
\newcommand{\GaugeH}{\mathcal C}
\newcommand{\HamCharge}{\mathfrak m_{\mathrm H}}
\newcommand{\HamChargeOmega}{\mathfrak m_{\mathrm H,\Omega}}
\newcommand{\PiOmega}{\Pi^{(\Omega)}}
\newcommand{\AY}{\mathcal A_Y}
\newcommand{\AZ}{\mathcal A_Z}
\newcommand{\Exterior}{\mathcal E_{\mathrm{ext}}}

\newtheorem{definition}{Definition}
\newtheorem{proposition}{Proposition}
\newtheorem{remark}{Remark}
\newtheorem{corollary}{Corollary}
\newtheorem{theorem}{Theorem}

\title{\textbf{The \TheoryName{}}\\[0.4em]
\large Inverse-Einstein-Response Quartic Closure Test}
\author{}
\date{}

\begin{document}

\maketitle

\begin{abstract}
The inverse-Einstein-response bridge-map test recorded a limited positive broader-screen verdict: the full currently recorded off-longitudinal \(Q/W\) remainder can be canceled by a solve-defined observer map, and the first surviving cubic longitudinal remainder can be removed by an explicit cubic completion. That note left two exact-closure burdens open. First, the solve itself had to be made honest as an admissible bridge map rather than as a purely observer-side imposition. Second, the first unresolved remainder
\[
\ReLongFour+\EinQuad[\SigmaIER]-8\pi G_N\,\DeltaTIER
\]
had to be computed explicitly enough to decide whether it vanishes, is only gauge/constraint exact, or remains negative.

The present note does both jobs. It fixes one observer-level admissible inverse-Einstein-response class: harmonic gauge, Coulomb-plus-regular right inverse, Hamiltonian-charge normalization for the monopole tail, and zero extra homogeneous radiative data. In particular, for any admissible source \(S_{\mu\nu}\) with Hamiltonian charge \(\HamCharge[S]\), the allowed inverse is required to split as
\[
\EinInv[S]=\SigmaC[S]+\SigmaReg[S],
\]
where \(\SigmaC[S]\) is the standard harmonic Coulomb monopole
\[
\SigmaC[S]_{00}=2U_S,\qquad
\SigmaC[S]_{0i}=0,\qquad
\SigmaC[S]_{ij}=2U_S\,\delta_{ij},
\qquad
U_S(r)=\frac{G_N\,\HamCharge[S]}{r},
\]
and where the regular part carries zero charge, decays one power faster, and contains no freely inserted homogeneous radiation.

With that class fixed, the unresolved observer tensor becomes
\[
\mathcal R_{\mu\nu}^{\mathrm{IER}}
=
\ReLongFour
+\mathcal Q_{\mu\nu}^{\mathrm C}
+\mathcal Q_{\mu\nu}^{\mathrm{mix}}
+\mathcal Q_{\mu\nu}^{\mathrm{reg}}
-8\pi G_N\,\DeltaTIER,
\]
where \(\mathcal Q^{\mathrm C}=\EinQuad[\SigmaC]\), \(\mathcal Q^{\mathrm{mix}}=2\EinBil(\SigmaC,\SigmaReg+\SigmaLongThree)\), and \(\mathcal Q^{\mathrm{reg}}=\EinQuad[\SigmaReg+\SigmaLongThree]\). Thus the first unresolved quartic burden is no longer an unspecified symbol: it is the sum of the Coulomb self-response, the mixed Coulomb-regular response, the purely regular/higher-order response, the quartic longitudinal remainder, and the matter-stress correction.

The decisive regime is the same purely rotational matter-free regime that forced the Coulomb tail in the first place:
\[
K\equiv0,\qquad
Q^I\equiv0,\qquad
\Omega_{ij}^T\not\equiv0,\qquad
T_{\mu\nu}^{\mathrm{matter}}=0.
\]
There the longitudinal pieces vanish, the matter correction vanishes, and the charge-free regular part decays like \(r^{-2}\). Hence the leading quartic remainder is the Coulomb self-response alone. Writing
\[
\HamChargeOmega\equiv \HamCharge[\ReOffTwo]
=
\frac{\kappa_\Omega}{4}\int_{\mathbb R^3}\abs{\Omega^T}^2\,d^3x>0,
\qquad
U_\Omega(r)=\frac{G_N\,\HamChargeOmega}{r},
\]
one finds on the exterior region of the source support that
\[
n^\mu n^\nu \mathcal R_{\mu\nu}^{\mathrm{IER}}
=
\frac{3\abs{\nabla U_\Omega}^2}{(1+2U_\Omega)^3}
+\mathcal O(r^{-5})
=
\frac{3G_N^2\HamChargeOmega^2}{r^4}
+\mathcal O(r^{-5})>0.
\]
Therefore the quartic remainder neither vanishes nor becomes only gauge/constraint exact in the admissible stationary class. The quartic-and-higher continuation is negative. The inverse-Einstein-response map remains the first broader observer-level bridge map to pass the old quadratic and cubic screens, but it still fails exact closure because the required Coulomb tail feeds a genuine quartic Einstein self-response back into the observer equation. The next live burden is therefore not a substrate-origin note for this map as if closure had succeeded; it is one layer deeper, namely a more fundamental nonlocal bridge-map redesign or substrate mechanism that can control or replace this self-response.
\end{abstract}

\section{Introduction}

The active sequence is now clear about what has and has not been settled. The positive quotient reduced-system, theorem, decay, and asymptotic line remains closed and is not reopened here. The bridge-first redesign, the matter-route-only redesign, and the first explicit local operational-metric candidate test also remain fixed in status. The inverse-Einstein-response bridge-map test then took the first genuinely broader step: it solved the linearized Einstein response against the full off-longitudinal observer remainder, fixed the necessary Coulomb tail by the Hamiltonian charge of the source, and added an explicit cubic longitudinal completion.

That broader-screen result was deliberately limited. It did not yet make the solve itself fully honest as an admissible bridge map. It also did not yet decide whether the first unresolved observer tensor
\[
\ReLongFour+\EinQuad[\SigmaIER]-8\pi G_N\,\DeltaTIER
\]
vanishes, is only gauge/constraint exact, or remains negative.

\paragraph{Framework and claim boundary.}
\TheoryProgram{} denotes the broader \Aether{} / \Aflow{} framework built on the claims that the \Aether{} is the underlying four-dimensional substrate of reality, the \Aflow{} is its intrinsic ordered motion, observed three-dimensional space is the local experiential slice of that deeper substrate, S-time is the experienced order of change arising from matter, light, and the \Aflow{}, and observed expansion is the three-dimensional appearance of deeper four-dimensional ordered motion. Within that framework, \TheoryName{} denotes the exact-closure theory in which the effective gravitational dynamics are adopted to be exactly Einsteinian with universal matter coupling. In the active sequence, \emph{adoption} means use of that established relativistic dynamics without claiming substrate derivation, while \emph{derivation} is reserved for a first-principles recovery from explicit substrate variables.

The present note stays strictly inside that boundary. It does not claim a completed substrate derivation. It does not reopen the positive quotient PDE line. It does not pretend that the inverse-Einstein-response solve already comes from explicit substrate microphysics. Its task is narrower and exact-closure driven: make the solve class observer-level admissible, compute the first unresolved quartic tensor honestly, and record the resulting verdict.

\section{Current Branch and the Admissible Solve Class}

\begin{definition}[Current redesigned bridge branch]
\label{def:branch}
The \emph{current redesigned bridge branch} is the branch fixed by the redesigned-bridge closure/tensor-verdict note, the matter-route redesign note, the broader operational-metric redesign note, and the inverse-Einstein-response bridge-map test:
\begin{enumerate}
    \item the quotient equations
    \begin{equation}
    \AY=0,
    \qquad
    \AZ=0,
    \label{eq:AYAZ}
    \end{equation}
    are imposed from the outset;
    \item the operational bridge metric is \(\ReBridge_{\mu\nu}\);
    \item the eliminated observer equation is
    \begin{equation}
    G_{\mu\nu}[\ReBridge]
    =
    8\pi G_N\,T_{\mu\nu}^{\mathrm{matter}}[\ReBridge,\psi]
    +\ReOff+\ReLongThree+\ReLongFour;
    \label{eq:oldobs}
    \end{equation}
    \item the inverse-Einstein-response correction is
    \begin{equation}
    \SigmaIER=\SigmaOff+\SigmaLongThree,
    \qquad
    \OpMetricIER=\ReBridge+\SigmaIER,
    \label{eq:iermetric}
    \end{equation}
    with
    \begin{equation}
    \EinLin(\SigmaOff)=-\ReOff,
    \qquad
    \EinLin(\SigmaLongThree)=-\ReLongThree.
    \label{eq:solveeqs}
    \end{equation}
\end{enumerate}
\end{definition}

\begin{definition}[Observer-admissible inverse-Einstein-response source]
\label{def:source}
A symmetric tensor \(S_{\mu\nu}\) on the current redesigned bridge branch is \emph{observer admissible for the inverse-Einstein-response solve} if:
\begin{enumerate}
    \item \(S_{\mu\nu}\) is symmetric and divergence controlled after the Einstein constraints and reduced slice equations are imposed;
    \item the Hamiltonian charge
    \begin{equation}
    \HamCharge[S]
    \equiv
    \int_{\Sigma_t} n^\mu n^\nu S_{\mu\nu}\,d\mu_\gamma
    \label{eq:charge}
    \end{equation}
    is finite;
    \item \(S_{\mu\nu}\) has one asymptotically flat observer chart in which the source admits a monopole coefficient determined by \(\HamCharge[S]\) and a charge-free remainder with one higher power of decay;
    \item no additional free radiative homogeneous data are inserted into the solve.
\end{enumerate}
\end{definition}

\begin{definition}[Admissible harmonic Coulomb-plus-regular right inverse]
\label{def:inverse}
Fix the harmonic gauge condition
\begin{equation}
\GaugeH_\nu(H)
\equiv
\nabla^\mu
\left(
H_{\mu\nu}
-\frac12 \ReBridge_{\mu\nu}(\ReBridge)^{\alpha\beta}H_{\alpha\beta}
\right)
=0.
\label{eq:gauge}
\end{equation}
For any observer-admissible source \(S_{\mu\nu}\), the \emph{admissible inverse-Einstein-response right inverse} is the specific solve
\begin{equation}
\EinInv[S]=\SigmaC[S]+\SigmaReg[S]
\label{eq:invsplit}
\end{equation}
with the following properties:
\begin{enumerate}
    \item \(\EinLin(\SigmaC[S]+\SigmaReg[S])=S\) and \(\GaugeH(\SigmaC[S]+\SigmaReg[S])=0\);
    \item in the asymptotic exact-flat chart, the Coulomb part is fixed by the Hamiltonian charge:
    \begin{equation}
    \SigmaC[S]_{00}=2U_S,
    \qquad
    \SigmaC[S]_{0i}=0,
    \qquad
    \SigmaC[S]_{ij}=2U_S\,\delta_{ij},
    \qquad
    U_S(r)=\frac{G_N\,\HamCharge[S]}{r};
    \label{eq:coulombtail}
    \end{equation}
    \item the regular part is charge free and decays one power faster:
    \begin{equation}
    \SigmaReg[S]=\mathcal O(r^{-2}),
    \qquad
    \partial\SigmaReg[S]=\mathcal O(r^{-3});
    \label{eq:regular}
    \end{equation}
    \item the homogeneous boundary data are fixed to zero beyond the forced Coulomb tail of \eqref{eq:coulombtail}.
\end{enumerate}
\end{definition}

\begin{remark}
Definition \ref{def:inverse} is the observer-level admissibility package missing from the previous note. The required Coulomb tail is not an optional observer-side decoration. Once \(\HamCharge[S]\neq0\), it is part of the admissible right-inverse class itself. The freedom that would otherwise survive in the homogeneous radiative sector is removed here by fixing that data to zero at the present exact-closure test.
\end{remark}

\begin{proposition}[The inverse-Einstein-response bridge map is now an admissible observer-level map]
\label{prop:adm}
At the present scope, the source tensors \(\ReOff\) and \(\ReLongThree\) are observer admissible in the sense of Definition \ref{def:source}. Hence the inverse-Einstein-response correction can be written as
\begin{equation}
\SigmaOff=\SigmaC[-\ReOff]+\SigmaReg[-\ReOff],
\qquad
\SigmaLongThree=\EinInv[-\ReLongThree],
\label{eq:iersplit}
\end{equation}
with the Coulomb coefficient of \(\SigmaOff\) fixed by \(\HamCharge[\ReOff]\) and with no extra free homogeneous data.
\end{proposition}

\begin{proof}
The previous inverse-Einstein-response note already established that \(\ReOff\) and \(\ReLongThree\) are the first active observer remainders on the current redesigned branch, that their Hamiltonian projections are well defined at the present scope, and that \(\ReOff\) carries the decisive nonzero Hamiltonian charge on the purely rotational regime. Definition \ref{def:inverse} then fixes the corresponding admissible right-inverse class uniquely at observer level: harmonic gauge, forced Coulomb monopole, charge-free faster-decaying regular part, and zero extra homogeneous radiation.
\end{proof}

\section{Exact Quartic Remainder Decomposition}

\begin{definition}[Polarized quadratic Einstein self-response]
\label{def:bil}
For any symmetric tensors \(H^{(1)}_{\mu\nu}\) and \(H^{(2)}_{\mu\nu}\), define the symmetric bilinear polarization of \(\EinQuad\) by
\begin{equation}
\EinBil(H^{(1)},H^{(2)})_{\mu\nu}
\equiv
\frac12
\left(
\EinQuad[H^{(1)}+H^{(2)}]_{\mu\nu}
-\EinQuad[H^{(1)}]_{\mu\nu}
-\EinQuad[H^{(2)}]_{\mu\nu}
\right).
\label{eq:bil}
\end{equation}
\end{definition}

\begin{definition}[Coulomb, mixed, and regular quartic pieces]
\label{def:qpieces}
Using the admissible decomposition \eqref{eq:iersplit}, define
\begin{align}
\mathcal Q_{\mu\nu}^{\mathrm C}
&\equiv
\EinQuad[\SigmaC[-\ReOff]]_{\mu\nu},
\label{eq:QC}
\\
\mathcal Q_{\mu\nu}^{\mathrm{mix}}
&\equiv
2\,\EinBil\!\left(\SigmaC[-\ReOff],\,\SigmaReg[-\ReOff]+\SigmaLongThree\right)_{\mu\nu},
\label{eq:Qmix}
\\
\mathcal Q_{\mu\nu}^{\mathrm{reg}}
&\equiv
\EinQuad[\SigmaReg[-\ReOff]+\SigmaLongThree]_{\mu\nu}.
\label{eq:Qreg}
\end{align}
\end{definition}

\begin{theorem}[The first unresolved tensor written explicitly]
\label{thm:decomp}
On reduced solutions of the current redesigned bridge branch,
\begin{equation}
\mathcal R_{\mu\nu}^{\mathrm{IER}}
\equiv
\ReLongFour+\EinQuad[\SigmaIER]_{\mu\nu}-8\pi G_N\,\DeltaTIER
\label{eq:defR}
\end{equation}
admits the explicit decomposition
\begin{equation}
\mathcal R_{\mu\nu}^{\mathrm{IER}}
=
\ReLongFour
+\mathcal Q_{\mu\nu}^{\mathrm C}
+\mathcal Q_{\mu\nu}^{\mathrm{mix}}
+\mathcal Q_{\mu\nu}^{\mathrm{reg}}
-8\pi G_N\,\DeltaTIER.
\label{eq:decomp}
\end{equation}
\end{theorem}

\begin{proof}
Insert the admissible split
\[
\SigmaIER=\SigmaC[-\ReOff]+\SigmaReg[-\ReOff]+\SigmaLongThree
\]
into \(\EinQuad[\SigmaIER]\). By Definition \ref{def:bil}, quadratic polarization gives exactly the sum of the pure Coulomb piece \eqref{eq:QC}, the mixed piece \eqref{eq:Qmix}, and the regular/higher-order piece \eqref{eq:Qreg}. Substituting that identity into \eqref{eq:defR} yields \eqref{eq:decomp}.
\end{proof}

\begin{remark}
Theorem \ref{thm:decomp} is the first main output of the present note. The previous manuscript only reduced the burden to the symbolic combination \(\ReLongFour+\EinQuad[\SigmaIER]-8\pi G_N\,\DeltaTIER\). The present note resolves that symbol into definite pieces, one of which is the Einstein self-response of the forced Coulomb tail itself.
\end{remark}

\section{Decisive Rotational-Regime Reduction}

\begin{definition}[Purely rotational matter-free regime]
\label{def:rot}
A reduced solution lies in the \emph{purely rotational matter-free regime} if
\begin{equation}
K\equiv0,
\qquad
Q^I\equiv0,
\qquad
\Omega_{ij}^T\not\equiv0,
\qquad
T_{\mu\nu}^{\mathrm{matter}}=0.
\label{eq:rot}
\end{equation}
\end{definition}

\begin{proposition}[What vanishes on the rotational regime]
\label{prop:rotzero}
On the regime of Definition \ref{def:rot},
\begin{equation}
\ReLongThree=0,
\qquad
\ReLongFour=0,
\qquad
\SigmaLongThree=0,
\qquad
\DeltaTIER=0.
\label{eq:rotzero}
\end{equation}
Moreover,
\begin{equation}
\HamCharge[\ReOffTwo]
=
\frac{\kappa_\Omega}{4}\int_{\mathbb R^3}\abs{\Omega^T}^2\,d^3x
>0.
\label{eq:rotcharge}
\end{equation}
\end{proposition}

\begin{proof}
This is exactly the convergence-free rotational regime isolated by the redesigned-bridge verdict, the first explicit operational-metric candidate test, and the inverse-Einstein-response bridge-map test. There the only active auxiliary source is the transverse vorticity block. The longitudinal trace/convergence variables vanish identically, so the cubic and quartic longitudinal remainders and their explicit cubic completion disappear. Matter is absent by definition, so \(\DeltaTIER=0\). The charge identity \eqref{eq:rotcharge} is the previously recorded Hamiltonian projection formula for the quadratic vorticity tensor.
\end{proof}

\begin{definition}[Rotational Coulomb tail]
\label{def:Uomega}
On the regime \eqref{eq:rot}, write
\begin{equation}
\HamChargeOmega
\equiv
\HamCharge[\ReOffTwo],
\qquad
U_\Omega(r)
\equiv
\frac{G_N\,\HamChargeOmega}{r}.
\label{eq:Uomega}
\end{equation}
Define the corresponding Coulomb solve by
\begin{equation}
\SigmaC_\Omega{}_{00}=2U_\Omega,
\qquad
\SigmaC_\Omega{}_{0i}=0,
\qquad
\SigmaC_\Omega{}_{ij}=2U_\Omega\,\delta_{ij}.
\label{eq:SigmaComega}
\end{equation}
\end{definition}

\begin{proposition}[Asymptotic order of the remaining pieces]
\label{prop:orders}
On the rotational regime, the admissible off-longitudinal solve splits as
\begin{equation}
\SigmaOff=\SigmaC_\Omega+\SigmaReg_\Omega,
\label{eq:offsplit}
\end{equation}
where
\begin{equation}
\SigmaReg_\Omega=\mathcal O(r^{-2}),
\qquad
\partial\SigmaReg_\Omega=\mathcal O(r^{-3}).
\label{eq:regomega}
\end{equation}
Consequently, on the exterior region \(\Exterior\) outside the source support,
\begin{equation}
\mathcal Q_{\mu\nu}^{\mathrm{mix}}=\mathcal O(r^{-5}),
\qquad
\mathcal Q_{\mu\nu}^{\mathrm{reg}}=\mathcal O(r^{-6}),
\label{eq:qorders}
\end{equation}
and therefore
\begin{equation}
\mathcal R_{\mu\nu}^{\mathrm{IER}}
=
\mathcal Q_{\mu\nu}^{\mathrm C}
+\mathcal O(r^{-5})
\qquad
\text{on }\Exterior.
\label{eq:lead}
\end{equation}
\end{proposition}

\begin{proof}
The monopole coefficient of \(\SigmaOff\) is fixed completely by \(\HamChargeOmega\), so the remainder after subtracting \(\SigmaC_\Omega\) is charge free and decays one power faster by Definition \ref{def:inverse}. Since \(\SigmaC_\Omega=\mathcal O(r^{-1})\) with \(\partial\SigmaC_\Omega=\mathcal O(r^{-2})\), while \(\SigmaReg_\Omega=\mathcal O(r^{-2})\) with \(\partial\SigmaReg_\Omega=\mathcal O(r^{-3})\), quadratic polarization gives \(\mathcal Q^{\mathrm{mix}}=\mathcal O(r^{-5})\) and \(\mathcal Q^{\mathrm{reg}}=\mathcal O(r^{-6})\). Proposition \ref{prop:rotzero} removes the longitudinal and matter terms, leaving \eqref{eq:lead}.
\end{proof}

\section{Explicit Coulomb Self-Response}

\begin{proposition}[Hamiltonian projection of the Coulomb self-response]
\label{prop:coulomb}
On the exterior region \(\Exterior\), the Hamiltonian projection of the pure Coulomb quartic term satisfies
\begin{equation}
n^\mu n^\nu \mathcal Q_{\mu\nu}^{\mathrm C}
=
-\frac{2\Delta U_\Omega}{(1+2U_\Omega)^2}
+\frac{3\abs{\nabla U_\Omega}^2}{(1+2U_\Omega)^3}.
\label{eq:coulomb}
\end{equation}
In particular, because \(\Delta U_\Omega=0\) on \(\Exterior\),
\begin{equation}
n^\mu n^\nu \mathcal Q_{\mu\nu}^{\mathrm C}
=
\frac{3\abs{\nabla U_\Omega}^2}{(1+2U_\Omega)^3}
=
\frac{3G_N^2\HamChargeOmega^2}{r^4}
+\mathcal O(r^{-5})
>0.
\label{eq:coulombpositive}
\end{equation}
\end{proposition}

\begin{proof}
On \(\Exterior\), the Coulomb solve \eqref{eq:SigmaComega} is static and shift free, so the metric \(\eta+\SigmaC_\Omega\) has
\[
K_{ij}=0,
\qquad
\gamma_{ij}=(1+2U_\Omega)\delta_{ij}.
\]
Hence its Hamiltonian projection is
\[
n^\mu n^\nu G_{\mu\nu}[\eta+\SigmaC_\Omega]
=
\frac12 R[\gamma].
\]
Write \(\gamma_{ij}=e^{2\omega}\delta_{ij}\) with \(\omega=\frac12\log(1+2U_\Omega)\). The scalar curvature of a conformally flat three-metric is
\[
R[\gamma]
=
e^{-2\omega}
\left(
-4\Delta\omega
-2\abs{\nabla\omega}^2
\right),
\]
which here becomes
\[
R[\gamma]
=
-\frac{4\Delta U_\Omega}{(1+2U_\Omega)^2}
+\frac{6\abs{\nabla U_\Omega}^2}{(1+2U_\Omega)^3}.
\]
Dividing by \(2\) gives \eqref{eq:coulomb}. On \(\Exterior\), \(U_\Omega\) is harmonic, so \(\Delta U_\Omega=0\). Because
\[
\abs{\nabla U_\Omega}^2
=
\frac{G_N^2\HamChargeOmega^2}{r^4},
\]
equation \eqref{eq:coulombpositive} follows. Finally, \(\EinLin(\SigmaC_\Omega)=0\) on \(\Exterior\) because the source is absent there, so this Hamiltonian projection is exactly the Hamiltonian projection of \(\mathcal Q^{\mathrm C}=\EinQuad[\SigmaC_\Omega]\).
\end{proof}

\begin{remark}
Proposition \ref{prop:coulomb} is the decisive quartic computation. The tail required to evade the old integrated Hamiltonian obstruction is not inert. Once the quadratic source has nonzero Hamiltonian charge, the corresponding Coulomb solve feeds back through the quadratic Einstein term with a positive \(r^{-4}\) Hamiltonian density.
\end{remark}

\section{Gauge/Constraint Exactness and the Quartic Verdict}

\begin{definition}[Admissible gauge/constraint exactness]
\label{def:pure}
We call a symmetric tensor \(R_{\mu\nu}\) \emph{admissibly gauge/constraint exact} if there exist an observer-admissible one-form \(Y_\mu\), the Einstein Hamiltonian and momentum constraints \(\mathcal H\), \(\mathcal M_\mu\), and the reduced slice equations \(E_{Q,I}=0\), \(E_\chi=0\), \(E_{W,i}=0\) such that
\begin{equation}
R_{\mu\nu}
=
\nabla_{(\mu}Y_{\nu)}
+\mathcal P_{\mu\nu}\mathcal H
+\mathcal P_{\mu\nu}{}^\rho \mathcal M_\rho
+\mathcal P_{\mu\nu}^{(Q)I}E_{Q,I}
+\mathcal P_{\mu\nu}^{(\chi)}E_\chi
+\mathcal P_{\mu\nu}^{(W)i}E_{W,i},
\label{eq:pure}
\end{equation}
for differential operators \(\mathcal P\) built from the reduced variables, with \(Y_\mu\) carrying no extra growing or radiative homogeneous data beyond the stationary asymptotic class fixed in Definition \ref{def:inverse}.
\end{definition}

\begin{theorem}[Quartic-and-higher inverse-Einstein-response verdict]
\label{thm:verdict}
The quartic-and-higher continuation of the inverse-Einstein-response bridge map is negative:
\begin{enumerate}
    \item the first unresolved remainder \(\mathcal R_{\mu\nu}^{\mathrm{IER}}\) does \emph{not} vanish on the purely rotational matter-free regime;
    \item \(\mathcal R_{\mu\nu}^{\mathrm{IER}}\) is not only admissibly gauge/constraint exact on that regime;
    \item therefore the solve-defined inverse-Einstein-response bridge map still does \emph{not} close the observer equation exactly.
\end{enumerate}
\end{theorem}

\begin{proof}
By Proposition \ref{prop:orders},
\[
\mathcal R_{\mu\nu}^{\mathrm{IER}}
=
\mathcal Q_{\mu\nu}^{\mathrm C}
+\mathcal O(r^{-5})
\qquad
\text{on }\Exterior
\]
along the purely rotational matter-free regime. Proposition \ref{prop:coulomb} then gives
\[
n^\mu n^\nu \mathcal R_{\mu\nu}^{\mathrm{IER}}
=
\frac{3G_N^2\HamChargeOmega^2}{r^4}
+\mathcal O(r^{-5})
>0
\]
for large \(r\), so \(\mathcal R_{\mu\nu}^{\mathrm{IER}}\not\equiv0\). This proves item (1).

For item (2), on the exterior region \(\Exterior\) the source support has been removed, so the Einstein constraints and reduced slice equations vanish there. Any admissibly gauge/constraint exact representation must therefore reduce on \(\Exterior\) to a pure gauge term \(\nabla_{(\mu}Y_{\nu)}\). In the stationary asymptotic class fixed by Definition \ref{def:inverse}, the allowed \(Y_\mu\) carries no growing time-dependent homogeneous mode, so on the static exterior regime one has
\[
n^\mu n^\nu \nabla_{(\mu}Y_{\nu)}=0.
\]
That cannot reproduce the positive Hamiltonian projection of \(\mathcal R_{\mu\nu}^{\mathrm{IER}}\) from Proposition \ref{prop:coulomb}. Hence \(\mathcal R_{\mu\nu}^{\mathrm{IER}}\) is not only admissibly gauge/constraint exact.

Item (3) is immediate: exact observer closure would require the first unresolved remainder to vanish or to collapse to an admissibly gauge/constraint exact form, and neither condition holds.
\end{proof}

\begin{corollary}[Status of the inverse-Einstein-response map]
\label{cor:status}
The inverse-Einstein-response map remains the first genuinely broader observer-level bridge map to survive the old quadratic off-longitudinal and cubic longitudinal screens, but it fails at exact-closure level once the required Coulomb tail is treated honestly inside the admissible solve class.
\end{corollary}

\begin{corollary}[What the next manuscript burden now is]
\label{cor:next}
Because the quartic continuation is negative, the next live manuscript burden is not a substrate-origin note for the present inverse-Einstein-response map as if that map had already closed positively. It is a deeper nonlocal bridge-map redesign, or an explicit substrate mechanism that changes the observer-level self-response itself rather than merely solving the same linear Einstein response with a forced Coulomb tail.
\end{corollary}

\section{Conclusion}

The previous inverse-Einstein-response bridge-map test was worth having because it changed the shape of the problem. It showed that the old decisive off-longitudinal obstruction was not the last word. A genuinely broader observer map could cancel the full recorded \(Q/W\) remainder and could remove the first surviving cubic longitudinal remainder by explicit completion. But that positive broader-screen result left one exact-closure question open: what does the first unresolved quartic tensor actually do once the solve is made honest?

The answer is now explicit. The solve can be made observer-level admissible by fixing harmonic gauge, choosing a Coulomb-plus-regular right inverse, determining the monopole tail by the Hamiltonian charge of the source, and forbidding any extra homogeneous radiation. With that class fixed, the unresolved observer tensor decomposes into the quartic longitudinal remainder, the Coulomb self-response, the mixed Coulomb-regular response, the purely regular higher-order response, and the matter-stress correction. On the same purely rotational matter-free regime that forced the Coulomb tail in the first place, the longitudinal and matter pieces vanish and the charge-free remainder decays faster, so the leading quartic term is exactly the Coulomb self-response.

That leading term is not harmless. Its Hamiltonian projection is positive,
\[
n^\mu n^\nu \mathcal R_{\mu\nu}^{\mathrm{IER}}
=
\frac{3G_N^2\HamChargeOmega^2}{r^4}
+\mathcal O(r^{-5}),
\]
so the quartic remainder neither vanishes nor reduces to an admissible gauge/constraint-exact expression. The quartic-and-higher continuation is therefore negative. The broader bridge-map idea was still the right next test, but it does not yet deliver exact closure. The present map remains observer-level and nonlocal, and its forced Coulomb tail generates its own genuine quartic Einstein self-response. The active exact-closure burden now moves one layer deeper again.

\end{document}
